\documentclass[11pt]{article}

% ---------- Packages ----------
\usepackage[utf8]{inputenc}
\usepackage[T1]{fontenc}
\usepackage{lmodern}
\usepackage[a4paper,margin=2.5cm]{geometry}
\usepackage{graphicx}
\usepackage{booktabs}
\usepackage{amsmath}
\usepackage{authblk}
\usepackage{caption}
\usepackage{float}
\usepackage{siunitx}
\usepackage[round]{natbib}
\usepackage{xcolor}
\usepackage[colorlinks=true,allcolors=blue]{hyperref}

\graphicspath{{../outputs/figures/}}

% Draft-note macro: highlights text the author still needs to write/refine.
\newcommand{\authornote}[1]{{\color{red}\textit{[#1]}}}

\captionsetup{font=small,labelfont=bf}

% ---------- Title ----------
\title{\textbf{Structurally Disadvantaged: A Statistical Analysis of
Batting Position and Team Qualification Bias in the IPL Orange Cap Award}}

\author[1]{Sandheep Sridar}
\affil[1]{\small Independent Researcher \\ \texttt{sandeepsridar@gmail.com}}

\date{\today}

\begin{document}

\maketitle

% ---------- Abstract ----------
\begin{abstract}
\noindent
The Orange Cap is the most sought-after individual batting award in the Indian
Premier League (IPL). In this paper we discuss how it systematically
disadvantages middle-order batsmen and high performers who do not qualify for
the play-offs. This bias stems from three advantages enjoyed by openers and
play-off teams: greater ball-faced volume, exclusive access to power-play
conditions, and additional match opportunities from play-off qualification.

Using ball-by-ball data from Cricsheet comprising 278{,}034 deliveries across
18 seasons, we show that power-play access and play-off matches are strong
determinants of the Orange Cap award. We are not comparing players or evaluating
player performance; we are assessing the fairness of the award's design.

The results confirm each hypothesis. Openers face a median of 80 more balls per
season than middle-order batsmen and score 58\% of their runs in the power-play
overs. When runs are normalised to a 14-match baseline, 77.6\% of the top batters
per season rank differently, confirming the advantage of match access. Players who
qualify for the play-offs play a median of two additional matches per season,
widening the margin for the Orange Cap. In all 18 seasons every Orange Cap winner
has batted in positions 1--3.

The findings demonstrate that match access, not batting skill alone, determines
the award outcome. The award structurally disincentivises middle-order batting
skill. Beyond the unfairness in individual recognition, this has broader
implications for the well-documented importance of the number-4 position at the
international level. We call for a re-evaluation of the Orange Cap criteria to
account for structural positional and match-count advantages.
\end{abstract}

\vspace{0.5em}
\noindent\textbf{Keywords:} cricket analytics; IPL; T20; batting position;
award fairness; sports analytics

\section{Introduction}

The Indian Premier League is the most popular T20 cricket league in the world
\citep{iplvaluation2026}. It was established in 2008 and has gained popularity over
time, with its total business value reaching \$18.5 billion in 2026, fuelled by a
\$6.4 billion media-rights deal and roughly \$600 million in advertising revenue
\citep{iplmediarights}. The league recognises individual performance with several
awards at the end of the season, and the most popular of these is the Orange Cap.

To popularise the T20 format, the IPL introduced several individual performance
awards to celebrate explosive batting. The Orange Cap is awarded to the batsman who
scores the most aggregate runs in the season. Introduced in the inaugural 2008
season, it is a single-criterion award --- most runs scored. There is no adjustment
for batting position, balls faced, or matches played. The leader gets the privilege
of wearing an orange cap during play, and with it comes significant media visibility
throughout the season. Ever since its inception, every winner has batted in positions
1--3, with 15 of 18 being openers.

The first six overs of the game are the power-play, where only two fielders are
allowed outside the circle, giving the batsman a large advantage in scoring runs.
Whether power-play access gives openers an advantage over middle-order batsmen, and
whether play-off qualification provides additional opportunities to score, are
questions this paper examines empirically.

Prior work on cricket analytics has examined batting position and performance from
multiple angles. \citet{pandey2023bayesian} used Bayesian survival analysis to study
middle-order batsman performance in ODI cricket. Batting role-aware performance
indices for T20 cricket have been proposed \citep{dppi2022}. Probabilistic approaches
to batting-order advantage have also been explored \citep{probabilistic2020}. However,
to the best of our knowledge, no prior study has examined whether a major T20 cricket
award is structurally biased by batting position and match access.

The remainder of this paper is organised as follows. Section~\ref{sec:methods}
describes the data source and methodology. Section~\ref{sec:results} presents results
across six analyses. Section~\ref{sec:discussion} discusses the implications of the
findings, including the role of the Super Striker award as a partial compensating
mechanism. Section~\ref{sec:conclusion} concludes with a call for re-evaluation of the
Orange Cap criteria.

\section{Data and Methodology}
\label{sec:methods}

\authornote{Draft to write. Cover: (i) Cricsheet data source ---
1{,}241 matches, 278{,}034 deliveries, seasons 2008--2025;
(ii) variable definitions --- batting position (rank of first appearance per
innings), position groups (Opener 1--2, Top Order 3, Middle Order 4--5,
Finisher 6+), phase (power-play overs 1--6, middle 7--15, death 16--20),
legal balls (wides excluded), play-off team definition;
(iii) qualification threshold (7+ matches per season);
(iv) statistical tests used --- chi-square, Kruskal--Wallis, Mann--Whitney~U,
Wilcoxon signed-rank, $p<0.05$.}

\subsection*{Data source}
We use publicly available ball-by-ball match data from Cricsheet \citep{cricsheet},
covering all IPL matches from the 2008 season through 2025 (18 complete seasons;
the ongoing 2026 season is excluded). After restricting to the first two innings of
each match and excluding super-overs, the dataset comprises 278{,}034 deliveries.

\subsection*{Variable definitions}
\begin{itemize}
  \item \textbf{Batting position} is the rank of a batter's first appearance in an
        innings. Per-season position is the average across the batter's innings.
  \item \textbf{Position group} maps average position to one of four groups:
        Opener (1--2), Top Order (3), Middle Order (4--5), and Finisher (6+).
  \item \textbf{Phase} is derived from the over number: power-play (overs 1--6),
        middle (7--15), and death (16--20).
  \item \textbf{Balls faced} counts only legal deliveries; wides are excluded, in
        line with the official IPL strike-rate convention.
  \item \textbf{Play-off team} is any team that appeared in a Qualifier, Eliminator,
        or Final in a given season.
\end{itemize}
A batter-season qualifies for population comparisons only if the batter played at
least seven matches that season, avoiding single-appearance noise.

\subsection*{Statistical methods}
We use the chi-square test for the winner-position distribution (Analysis~A);
Kruskal--Wallis with Mann--Whitney~U post-hoc tests for differences across position
groups (Analyses~B, C); the Wilcoxon signed-rank test for paired actual-versus-
normalised rankings (Analysis~D); and one-sided Mann--Whitney~U tests for match-count
and run differences by play-off status (Analyses~E, F). We adopt a significance
threshold of $p<0.05$.

\section{Results}
\label{sec:results}

\authornote{Draft prose to write around each figure/table. The figures, tables and
key numbers below are final from the analysis pipeline.}

\subsection{Batting position of Orange Cap winners}
\label{sec:resA}
In all 18 seasons, every Orange Cap winner batted in positions 1--3; 15 of 18 were
openers (positions 1--2) and the remaining three batted at position~3
(Figure~\ref{fig:winners}). No middle-order or finishing batsman has ever won. A
chi-square test confirms the winner-position distribution differs significantly from
the population distribution of qualified run-scorers
($\chi^2 = 49.61$, $p < 0.001$, $df = 3$).

\begin{figure}[H]
  \centering
  \includegraphics[width=0.85\linewidth]{fig1_winner_positions.png}
  \caption{Distribution of Orange Cap winners by batting position group,
  2008--2025. Every winner batted in positions 1--3.}
  \label{fig:winners}
\end{figure}

\subsection{Balls faced across position groups}
\label{sec:resB}
Openers face a median of 284 legal balls per season, versus 204 for middle-order
batsmen --- a gap of 80 balls (Table~\ref{tab:balls}, Figure~\ref{fig:balls}). A
Kruskal--Wallis test confirms the difference across groups is significant
($H = 1051.8$, $p < 0.001$, $n = 2783$).

\begin{table}[H]
  \centering
  \caption{Balls faced per season by position group (legal balls; 7+ match
  qualified batter-seasons, 2008--2025).}
  \label{tab:balls}
  \begin{tabular}{lrrrr}
    \toprule
    Position group & Batter-seasons & Median balls & Mean balls & Mean runs \\
    \midrule
    Opener (1--2)      & 224 & 284 & 286.1 & 386.2 \\
    Top Order (3)      & 183 & 235 & 250.7 & 339.3 \\
    Middle Order (4--5)& 372 & 204 & 209.7 & 282.4 \\
    Finisher (6+)      & 376 & 101 & 112.5 & 155.8 \\
    \bottomrule
  \end{tabular}
\end{table}

\begin{figure}[H]
  \centering
  \includegraphics[width=0.85\linewidth]{fig2_balls_faced.png}
  \caption{Distribution of balls faced per season by position group.}
  \label{fig:balls}
\end{figure}

\subsection{Power-play run concentration}
\label{sec:resC}
Openers score 57.5\% of their seasonal runs in the power-play, compared with just
16.2\% for middle-order batsmen and 2.4\% for finishers
(Table~\ref{tab:powerplay}, Figure~\ref{fig:powerplay}). The difference in power-play
run share across groups is significant (Kruskal--Wallis, $p < 0.001$).

\begin{table}[H]
  \centering
  \caption{Share of seasonal runs by phase and position group (\%).}
  \label{tab:powerplay}
  \begin{tabular}{lrrr}
    \toprule
    Position group & Power-play & Middle & Death \\
    \midrule
    Opener (1--2)      & 57.5 & 35.4 & 7.1 \\
    Top Order (3)      & 45.1 & 44.3 & 10.6 \\
    Middle Order (4--5)& 16.2 & 56.4 & 27.3 \\
    Finisher (6+)      & 2.4 & 34.6 & 63.0 \\
    \bottomrule
  \end{tabular}
\end{table}

\begin{figure}[H]
  \centering
  \includegraphics[width=0.85\linewidth]{fig3_powerplay_concentration.png}
  \caption{Runs by phase as a share of seasonal total, by position group.}
  \label{fig:powerplay}
\end{figure}

\subsection{Runs versus batting position}
\label{sec:resD1}
Across all qualified batter-seasons, seasonal runs correlate negatively with average
batting position ($r = -0.61$, $p < 0.001$), as shown in Figure~\ref{fig:scatter}.

\begin{figure}[H]
  \centering
  \includegraphics[width=0.85\linewidth]{fig4_runs_vs_position.png}
  \caption{Seasonal runs versus average batting position, all qualified
  batter-seasons. Negative correlation ($r = -0.61$).}
  \label{fig:scatter}
\end{figure}

\subsection{Normalised-ranking counterfactual}
\label{sec:resD2}
Restricting to league-stage matches and normalising each batter's runs to a common
14-match baseline ($\text{norm\_runs} = \text{league\_runs} \times 14 /
\text{league\_matches}$), 77.6\% of the top-10 rankings per season change order. A
Wilcoxon signed-rank test confirms actual and normalised ranks differ significantly
($W = 1914.5$, $p < 0.001$, $n = 183$). Figure~\ref{fig:normalised} shows the
ranking shifts; 95\% confidence intervals account for batters who played fewer than
14 matches.

\begin{figure}[H]
  \centering
  \includegraphics[width=0.85\linewidth]{fig5_normalised_rankings.png}
  \caption{Actual versus normalised rankings of top run-scorers after normalising to
  a 14-match baseline.}
  \label{fig:normalised}
\end{figure}

\subsection{Play-off match advantage}
\label{sec:resE}
Play-off teams play a median of two additional matches per season relative to
non-play-off teams (Mann--Whitney~U, $p < 0.001$; $n_{\text{playoff}} = 72$,
$n_{\text{non}} = 84$), as shown in Figure~\ref{fig:matchcount}. Given a typical
opener pace of roughly 40 runs per match, two extra matches translate to a
material run cushion in a close Orange Cap race.

\begin{figure}[H]
  \centering
  \includegraphics[width=0.85\linewidth]{fig6_match_count.png}
  \caption{Matches played per season by play-off and non-play-off teams.}
  \label{fig:matchcount}
\end{figure}

\subsection{Non-play-off elite batsmen}
\label{sec:resF}
We identify 23 batter-seasons in which a top-5 run-scorer played for a non-play-off
team and was therefore capped at the 14-match league stage. Projecting their runs to
a 16-match schedule (the median for play-off teams) shows several would have
challenged or exceeded the actual Orange Cap total. Notably, a one-sided
Mann--Whitney~U test does \emph{not} find that play-off top-5 batters out-score
non-play-off top-5 batters ($p = 0.95$; medians 552 vs.\ 625): elite non-play-off
batsmen are not less productive --- they simply have fewer matches in which to
accumulate runs.

\section{Discussion}
\label{sec:discussion}

\authornote{Draft to write. Cover: (i) implications for award design --- the award
rewards volume/access over skill; (ii) the Super Striker award as a partial but
inadequate compensating mechanism (introduced 2018, $\geq$100-ball threshold, last
three winners 2024--2026 were openers --- converging toward the same pattern);
(iii) the broader importance of the number-4 position at international level;
(iv) limitations --- Impact Player rule (2023+), DLS-shortened matches, opposition
quality not modelled, batting position volatility.}

\subsection*{The Super Striker award}
The IPL introduced the Super Striker award in 2018, given to the batsman with the
highest strike rate (minimum 100 balls faced). Its existence is an implicit
acknowledgement that pure run accumulation does not capture batting quality. However,
it does not adequately compensate for the Orange Cap bias: it is a secondary award
with far less prestige, the 100-ball threshold still excludes many middle-order
batsmen, and the last three winners (2024--2026) have all been openers, converging
toward the same positional pattern as the Orange Cap itself.

\section{Conclusion}
\label{sec:conclusion}

\authornote{Draft to write. Summarise the structural bias, restate that match access
(position + play-off qualification) --- not skill alone --- determines the award, and
call for a re-evaluation of the Orange Cap criteria. Tease Paper~2 (a bias-corrected
Best Batsman Index).}

% ---------- References ----------
\bibliographystyle{plainnat}
\bibliography{references}

\end{document}
